\documentclass[11pt,a4paper]{article}

% =========================
% Packages
% =========================
\usepackage[utf8]{inputenc}
\usepackage[T1]{fontenc}
\usepackage[french]{babel}
\usepackage{geometry}
\usepackage{array}
\usepackage{longtable}
\usepackage{hyperref}
\usepackage[table]{xcolor}
\usepackage{listings}
\usepackage{enumitem}

\geometry{margin=2.2cm}

% =========================
% Couleurs
% =========================
\definecolor{straussblue}{RGB}{36,63,196}
\definecolor{lightblue}{RGB}{234,245,255}
\definecolor{lightgray}{RGB}{245,245,245}
\definecolor{codegray}{RGB}{248,248,248}

% =========================
% Style code
% =========================
\lstset{
    backgroundcolor=\color{codegray},
    basicstyle=\ttfamily\small,
    frame=single,
    breaklines=true,
    columns=fullflexible,
    keepspaces=true,
    showstringspaces=false
}

% =========================
% Infos document
% =========================
\title{\textbf{Documentation du projet ARIA ODM Builder}}
\author{Institut Strauss}
\date{\today}

\begin{document}

\maketitle

\tableofcontents
\newpage

% ============================================================
\section{Objectif du projet}
% ============================================================

\textbf{ARIA ODM Builder} est une application Streamlit permettant de construire
des exports patients structurés à partir de fichiers issus d'ARIA/MOSAIQ.

L'application permet notamment de :

\begin{itemize}
    \item importer un fichier \texttt{traitement\_patient} ;
    \item importer un fichier \texttt{formulaire\_patient} ;
    \item sélectionner une cohorte à partir de codes CIM10 ;
    \item utiliser un fichier de configuration \texttt{mapping.csv} ;
    \item sélectionner les colonnes formulaire utiles ;
    \item générer des variables temporelles avant, pendant et après radiothérapie ;
    \item produire des variables cumulées ;
    \item contrôler la qualité des jointures et des dates ;
    \item exporter un fichier final patient ;
    \item exporter un rapport de preuve ;
    \item sauvegarder ou réimporter un profil JSON.
\end{itemize}

% ============================================================
\section{Structure du projet}
% ============================================================

La structure recommandée du projet est la suivante :

\begin{lstlisting}
CODE/
├── app.py
├── requirements.txt
├── SCENARIOS.md
├── SCENARIOS.tex
├── inputs/
│   ├── traitement_patient.csv
│   └── formulaire_patient.csv
├── conf/
│   ├── mapping.csv
│   └── profiles/
│       └── profil_exemple.json
├── pictures/
│   └── logo_strauss.png
├── assets/
└── utils/
    ├── __init__.py
    ├── text.py
    ├── clean.py
    ├── load.py
    ├── mapping.py
    ├── cohort.py
    ├── temporal.py
    ├── quality.py
    ├── export.py
    ├── profile.py
    └── display.py
\end{lstlisting}

% ============================================================
\section{Rôle des principaux fichiers}
% ============================================================

\subsection{\texttt{app.py}}

Le fichier \texttt{app.py} est le fichier principal de l'application Streamlit.

Il contient principalement :

\begin{itemize}
    \item la configuration de l'interface ;
    \item les onglets de navigation ;
    \item les boutons ;
    \item les zones d'import ;
    \item l'enchaînement général du workflow.
\end{itemize}

L'application se lance avec :

\begin{lstlisting}
streamlit run app.py
\end{lstlisting}

\subsection{\texttt{requirements.txt}}

Le fichier \texttt{requirements.txt} contient les librairies nécessaires au projet.

Exemple :

\begin{lstlisting}
streamlit
pandas
openpyxl
xlsxwriter
skrub
black
\end{lstlisting}

Installation :

\begin{lstlisting}
pip install -r requirements.txt
\end{lstlisting}

\subsection{\texttt{inputs/}}

Le dossier \texttt{inputs/} est destiné aux fichiers d'entrée patients.

Exemple :

\begin{lstlisting}
inputs/
├── traitement_patient.csv
└── formulaire_patient.csv
\end{lstlisting}

Ces fichiers peuvent ensuite être chargés depuis l'onglet \textbf{Import} de
l'application.

\subsection{\texttt{conf/}}

Le dossier \texttt{conf/} est destiné aux fichiers de configuration.

Exemple :

\begin{lstlisting}
conf/
├── mapping.csv
└── profiles/
    └── profil_exemple.json
\end{lstlisting}

Le fichier \texttt{mapping.csv} remplace l'ancien fichier Excel de règles. Il
est plus simple à versionner, plus léger et plus compatible entre systèmes.

\subsection{\texttt{pictures/}}

Le dossier \texttt{pictures/} contient le logo ou les images utilisées par
l'application.

Exemple :

\begin{lstlisting}
pictures/
└── logo_strauss.png
\end{lstlisting}

Le logo est chargé avec \texttt{st.image()} depuis ce dossier.

\subsection{\texttt{utils/}}

Le dossier \texttt{utils/} contient les fonctions utilitaires du projet.

Le fichier \texttt{\_\_init\_\_.py} permet à Python de reconnaître
\texttt{utils/} comme un package.

\begin{longtable}{|p{0.30\textwidth}|p{0.62\textwidth}|}
\hline
\rowcolor{lightblue}
\textbf{Fichier} & \textbf{Rôle} \\
\hline
\endfirsthead

\hline
\rowcolor{lightblue}
\textbf{Fichier} & \textbf{Rôle} \\
\hline
\endhead

\texttt{utils/text.py}
&
Fonctions de normalisation texte, gestion des accents, correction du mojibake,
recherche flexible de colonnes.
\\
\hline

\texttt{utils/clean.py}
&
Nettoyage des données, normalisation des clés patient, normalisation des codes
CIM10, masques de valeurs non vides, nettoyage contrôlé avec skrub.
\\
\hline

\texttt{utils/load.py}
&
Lecture des fichiers CSV, Excel et ZIP, lecture mémoire-sûre, lecture streaming
des gros fichiers.
\\
\hline

\texttt{utils/mapping.py}
&
Gestion du \texttt{mapping.csv}, aide à la sélection CIM10, suggestions de
colonnes, pertinence des colonnes formulaire.
\\
\hline

\texttt{utils/cohort.py}
&
Détection des colonnes du fichier traitement, filtrage CIM10, construction de la
base patient.
\\
\hline

\texttt{utils/temporal.py}
&
Gestion des fenêtres temporelles, construction du formulaire long, génération
des variables avant RT, pendant RT, après RT et cumulées.
\\
\hline

\texttt{utils/quality.py}
&
Création du rapport qualité et du pipeline de preuve.
\\
\hline

\texttt{utils/export.py}
&
Exports CSV et Excel, fichier final patient, rapport de preuve.
\\
\hline

\texttt{utils/profile.py}
&
Sauvegarde et chargement des profils JSON.
\\
\hline

\texttt{utils/display.py}
&
Configuration visuelle, CSS léger, logo et en-tête de l'application.
\\
\hline

\end{longtable}

% ============================================================
\section{Installation}
% ============================================================

\subsection{Se placer dans le dossier du projet}

Dans un terminal :

\begin{lstlisting}
cd chemin/vers/ModifiedV1
\end{lstlisting}

\subsection{Installer les dépendances}

Commande standard :

\begin{lstlisting}
pip install -r requirements.txt
\end{lstlisting}

Si la commande \texttt{pip} n'est pas reconnue :

\begin{lstlisting}
python -m pip install -r requirements.txt
\end{lstlisting}

Sous Windows avec un chemin Python explicite :

\begin{lstlisting}
& "C:\Program Files\Python314\python.exe" -m pip install -r requirements.txt
\end{lstlisting}

% ============================================================
\section{Lancement de l'application}
% ============================================================

Commande standard :

\begin{lstlisting}
streamlit run app.py
\end{lstlisting}

Commande de secours Windows :

\begin{lstlisting}
& "C:\Program Files\Python314\python.exe" -m streamlit run app.py
\end{lstlisting}

Une fois lancée, l'application s'ouvre dans le navigateur.

% ============================================================
\section{Organisation de l'interface}
% ============================================================

L'application est organisée en plusieurs onglets :

\begin{center}
\textbf{Accueil}
$\rightarrow$
\textbf{Import}
$\rightarrow$
\textbf{Construction}
$\rightarrow$
\textbf{Contrôle qualité}
$\rightarrow$
\textbf{Sources}
$\rightarrow$
\textbf{Profil}
\end{center}

\subsection{Accueil}

L'onglet \textbf{Accueil} présente le contexte général du projet et le déroulement
du workflow.

\subsection{Import}

L'onglet \textbf{Import} permet de charger :

\begin{itemize}
    \item le fichier \texttt{traitement\_patient} ;
    \item le fichier \texttt{formulaire\_patient} ;
    \item le fichier \texttt{mapping.csv} ;
    \item éventuellement un profil JSON.
\end{itemize}

\subsection{Construction}

L'onglet \textbf{Construction} permet de :

\begin{itemize}
    \item choisir le mode CIM10 ;
    \item filtrer la cohorte ;
    \item sélectionner les colonnes formulaire ;
    \item définir les temporalités ;
    \item lancer le calcul avec le bouton \textbf{Construire / recalculer}.
\end{itemize}

Le bouton \textbf{Construire / recalculer} évite que l'application relance tout le
pipeline à chaque modification d'un widget.

\subsection{Contrôle qualité}

L'onglet \textbf{Contrôle qualité} affiche les vérifications principales :

\begin{itemize}
    \item patients traitement retrouvés dans le formulaire ;
    \item patients cohorte retrouvés dans le formulaire ;
    \item dates de début et de fin de radiothérapie ;
    \item valeurs formulaire conservées ;
    \item doublons ;
    \item dates invalides.
\end{itemize}

\subsection{Sources}

L'onglet \textbf{Sources} permet de documenter les colonnes générées et leur
origine.

\subsection{Profil}

L'onglet \textbf{Profil} permet de sauvegarder ou réimporter une configuration au
format JSON.

% ============================================================
\section{Fichiers attendus}
% ============================================================

\subsection{Fichier \texttt{traitement\_patient}}

Le fichier \texttt{traitement\_patient} doit contenir les informations de
traitement.

Colonnes importantes attendues ou détectées automatiquement :

\begin{itemize}
    \item identifiant patient ;
    \item code CIM10 ;
    \item date de première fraction ;
    \item date de dernière fraction ;
    \item dose ;
    \item nombre de fractions.
\end{itemize}

\subsection{Fichier \texttt{formulaire\_patient}}

Le fichier \texttt{formulaire\_patient} doit contenir les formulaires patients.

Colonnes importantes attendues :

\begin{itemize}
    \item \texttt{pt\_id} ;
    \item \texttt{date\_event} ;
    \item colonnes cliniques ou formulaire à exploiter.
\end{itemize}

\subsection{Fichier \texttt{mapping.csv}}

Le fichier \texttt{mapping.csv} sert à aider la sélection des codes CIM10, des
localisations et des colonnes formulaire pertinentes.

Il est placé de préférence dans :

\begin{lstlisting}
conf/mapping.csv
\end{lstlisting}

Format recommandé :

\begin{lstlisting}
CIM10;Description;Localisation;Requete;Phase;Commentaire
C61;Tumeur maligne de la prostate;Prostate;PSA;Cumul;
C61;Tumeur maligne de la prostate;Prostate;Dysurie;Aigu;
\end{lstlisting}

% ============================================================
\section{Profils JSON}
% ============================================================

Les profils JSON servent à sauvegarder une configuration d'extraction.

Ils peuvent contenir :

\begin{itemize}
    \item les colonnes sélectionnées ;
    \item les noms d'export ;
    \item les temporalités choisies ;
    \item les paramètres de fenêtre ;
    \item les choix de construction.
\end{itemize}

Ils sont placés de préférence dans :

\begin{lstlisting}
conf/profiles/
\end{lstlisting}

Exemple :

\begin{lstlisting}
conf/profiles/profil_prostate.json
\end{lstlisting}

% ============================================================
\section{Formatage du code avec Black}
% ============================================================

Le projet utilise Black pour formater automatiquement le code Python.

Installation :

\begin{lstlisting}
pip install black
\end{lstlisting}

Formatage de tout le projet :

\begin{lstlisting}
python -m black app.py utils
\end{lstlisting}

Sous Windows avec un chemin Python explicite :

\begin{lstlisting}
& "C:\Program Files\Python314\python.exe" -m black app.py utils
\end{lstlisting}

Si Black affiche :

\begin{lstlisting}
reformatted app.py
\end{lstlisting}

ou :

\begin{lstlisting}
All done!
\end{lstlisting}

le formatage est terminé.

% ============================================================
\section{Scénarios utilisateur}
% ============================================================

Le fichier \texttt{SCENARIOS.md} décrit les parcours utilisateurs sous forme de
tableau :

\begin{center}
\textbf{Narration utilisateur} \quad | \quad \textbf{Action dans l'application}
\end{center}

Ces scénarios servent à justifier la structure de l'application.

\begin{longtable}{|p{0.45\textwidth}|p{0.45\textwidth}|}
\hline
\rowcolor{lightblue}
\textbf{Narration utilisateur} & \textbf{Action dans l'application} \\
\hline
\endfirsthead

\hline
\rowcolor{lightblue}
\textbf{Narration utilisateur} & \textbf{Action dans l'application} \\
\hline
\endhead

L'utilisateur arrive sur l'application et veut comprendre l'objectif général.
&
Il consulte l'onglet \textbf{Accueil}, qui présente le workflow global.
\\
\hline

L'utilisateur doit importer les données patients.
&
Il va dans l'onglet \textbf{Import} et charge les fichiers traitement et formulaire.
\\
\hline

L'utilisateur veut construire une cohorte CIM10.
&
Il va dans \textbf{Construction}, saisit le ou les codes CIM10, puis clique sur
\textbf{Construire / recalculer}.
\\
\hline

L'utilisateur veut vérifier les données avant export.
&
Il consulte l'onglet \textbf{Contrôle qualité}.
\\
\hline

L'utilisateur veut comprendre l'origine des colonnes générées.
&
Il consulte l'onglet \textbf{Sources}.
\\
\hline

L'utilisateur veut réutiliser la configuration plus tard.
&
Il exporte ou importe un profil JSON depuis l'onglet \textbf{Profil}.
\\
\hline

\end{longtable}

Ces scénarios permettent de vérifier que l'application suit un parcours logique
pour l'utilisateur.





\end{document}
